\begin{textsample}{POS Dim 3 – rs – Score 15.00 – rs/rs\_inf\_4.txt}  \label{ex:f3_pos_249}
6.2 \textbf{Crianças} e \textbf{Infâncias} As \textbf{crianças} são sujeitos históricos, de direitos e \textbf{desejos}, que vivem e se desenvolvem nos contextos sociais e culturais em que estão inseridas.
Nessas condições, fazem amizades, \textbf{brincam}, desejam, aprendem, observam, experimentam, questionam, constroem sentidos sobre o mundo e sobre suas identidades pessoais e coletivas, produzindo cultura.
As \textbf{crianças} utilizam diversas linguagens para construir conhecimentos e buscam compreender o mundo através das relações e interações que estabelecem com os \textbf{adultos} e com outras \textbf{crianças} de diferentes idades, da mesma forma com o ambiente. ( BRASIL, 2009 ).
Os \textbf{bebês}, as \textbf{crianças} bem \textbf{pequenas} e as \textbf{crianças} \textbf{pequenas} são sujeitos que necessitam de atenção, proteção, alimentação, \textbf{brincadeiras}, \textbf{higiene}, \textbf{escuta} e \textbf{afeto}.
O fato de serem simultaneamente frágeis e potentes em relação ao mundo, de serem biologicamente sociais, as torna reféns da interação, da presença efetiva do outro e, principalmente, do investimento afetivo dado pela confiança do outro. ( BRASIL, 2009, p. 23 ).
As \textbf{crianças} são seres criativos e ativos e vivem suas \textbf{infâncias} no presente, não se resumindo a serem preparadas para o futuro.
Através das interações e da \textbf{brincadeira}, as aprendizagens e o desenvolvimento se constituem e se ampliam.
Não há um modo padronizado e único de viver a \textbf{infância}, por isso compreende -se que há diversas \textbf{infâncias}, assim como são diversas as realidades culturais, sociais, econômicas e políticas da sociedade em que se inserem.
Portanto, > [... ] temos concebido as \textbf{crianças} como \textbf{seres} humanos concretos, um corpo presente no aqui e agora em interação com outros, portanto, com direitos civis.
As \textbf{infâncias}, temos pensado como a forma específica de conceber, produzir e legitimar as experiências das \textbf{crianças}.
Assim, falamos em \textbf{infâncias} no plural, pois elas são vividas de modo muito diverso.
Ser \textbf{criança} não implica em ter que vivenciar um único tipo de \textbf{infância}.
As \textbf{crianças}, por serem \textbf{crianças}, não estão condicionadas as mesmas experiências. ( BRASIL, 2009, p.22 ).
A \textbf{infância} não é vista apenas como uma etapa da vida ou um momento do desenvolvimento das pessoas em uma determinada faixa \textbf{etária}, que precisa ser superada e se encerra com a juventude.
Ela deixa \textbf{marcas} que permanecem ao longo da vida e constituem os seres humanos, marcando um jeito de ser e estar no mundo.
Como experiência, a \textbf{infância} não é vivida da mesma maneira por todas as \textbf{crianças}, já que elas são diferentes entre si e vivem em contextos sociais e culturais diferentes e são marcadas pelo pertencimento de classe social, etnia, gênero.
Uma concepção de \textbf{infância} plural, que percebe as \textbf{crianças} como sujeitos ativos, que participam e intervêm no meio, entende que através de suas ações as \textbf{crianças} reelaboram, recriam e agem sobre o mundo e que seus processos de interação envolvem o criar e o transformar.
Pela \textbf{brincadeira}, as \textbf{crianças} incorporam os elementos do mundo em que vivem, ao mesmo tempo em que agem sobre eles e estabelecem relações sociais e aprendizagens.
Na elaboração da proposta pedagógica, é importante destacar a diferença entre \textbf{crianças} e alunos, uma vez que frequentemente essas palavras são compreendidas como sinônimas.
Muitas vezes, ao ingressarem na escola, as \textbf{crianças} passam a ser vistas apenas como alunos, porém é preciso ter clareza que as \textbf{crianças} têm direitos a serem garantidos, que são os de se desenvolver e aprender por meio da experiência peculiar de viverem suas \textbf{infâncias}.
Nessa perspectiva, o Referencial Curricular Gaúcho compreende que a \textbf{criança} é o centro do planejamento curricular, sujeito de direitos que se desenvolve nas interações, relações e práticas cotidianas, com singularidades próprias.
O \textbf{brincar}, como linguagem própria da \textbf{infância}, assim como o \textbf{cuidado} e as experiências diversas com os saberes dos diferentes campos, oportunizam o desenvolvimento integral e saudável das \textbf{crianças}.

% matched lemmas: adulto, afeto, bebê, brincadeira, brincar, criança, cuidado, desejo, escuta, etário, higiene, infância, marca, pequeno, seres
\end{textsample}
